\documentclass[10pt,a4paper]{amsart}
\usepackage[margin=19mm]{geometry}
\usepackage{amsmath,amssymb,mathtools}
\usepackage[scr=rsfs,cal=euler]{mathalfa}
\usepackage{xfrac}
\usepackage[unicode,hidelinks,bookmarks=false]{hyperref}
\newcommand{\ie}{i.e.\ }
\newcommand{\cf}{cf.\ }
\newcommand{\resp}{respectively\ }
\newcommand{\Subsection}{\subsection}
\providecommand{\nobreakdash}{\nobreak\hbox{-}}
\setlength{\emergencystretch}{2em}
\multlinegap0pt
\usepackage{bm}
\usepackage{bbm}
\usepackage{dsfont}
\usepackage{xspace}
\usepackage{cancel}
\usepackage{mathtools}
\usepackage{amsthm}
\makeatletter
\@ifundefined{theorem}{\newtheorem{theorem}{Theorem}}{}
\@ifundefined{corollary}{\newtheorem{corollary}[theorem]{Corollary}}{}
\@ifundefined{lemma}{\newtheorem{lemma}[theorem]{Lemma}}{}
\@ifundefined{proposition}{\newtheorem{proposition}[theorem]{Proposition}}{}
\makeatother
\title[Asymptotic stability of symmetric flows with viscous inflow ]{Asymptotic stability of symmetric flows with viscous inflow boundary condition}
\author{Yan Guo, Zhuolun Yang}
\date{Source: arXiv:2602.17059v2 (CC BY 4.0)}
\begin{document}
\maketitle

\noindent\emph{Source and licence.} Asymptotic stability of symmetric flows with viscous inflow boundary condition. Version arXiv:2602.17059v2 (CC BY 4.0), distributed under the Creative Commons Attribution 4.0 International licence, as recorded in the arXiv licence metadata for this exact version and shown on the abstract page https://arxiv.org/abs/2602.17059v2. Full rights notice: licenses/arxiv-2602.17059-CC-BY-4.0.txt.\par\smallskip

\noindent\textit{Editorial scope.} Counted unit: the Proposition whose source label is \texttt{prop:phi\_t} in the article (source0.tex lines 1664--1669, source label \texttt{prop:phi\_t}), delivered together with the article's own complete original proof of it (source0.tex lines 1670--1759, one \texttt{proof} environment of 90 lines). No mathematics is written, repaired or re-derived: statement and proof are the article's own text line for line. Required context retained uncounted: Linearized\_NS (lines 232--237); Navier\_bc2 (lines 250--252); ss\_estimates\_H (lines 360--367); H1\_control\_by\_R (lines 1040--1042); est:H\_weighted\_Hardy (lines 1077--1085); est:mix\_norm\_H (lines 1090--1095); est:mix\_norm\_Hardy (lines 1103--1108); est:mix\_norm\_misc (lines 1277--1284); est:g\_basic (lines 1613--1615); est:g\_x (lines 1618--1622). Every definition, display and label of those lines is retained unchanged. Omitted from this package and not redistributed: 1--231, 238--249, 253--359, 368--1039, 1043--1076, 1086--1089, 1096--1102, 1109--1276, 1285--1612, 1616--1617, 1623--1663, 1760--2850. Nothing inside a retained excerpt is omitted or altered: every retained body is delivered exactly as the article's lines are written, so no omission receipt of any kind accompanies this package. Citations: the retained excerpts carry no citation command at all, so no bibliography entry of the article is transcribed and no reference is redistributed. Reference adaptation: none was needed and none was made; every reference inside the delivered unit resolves inside it, so no unresolved reference or citation is produced. Printed slips kept literal: no printed word, symbol, hypothesis or number of the counted statement or of its proof was altered. Layout and class: the article's own class is \texttt{amsart} with packages including amssymb, mathrsfs, color, esint, hyperref, verbatim, cancel, geometry, tikz, cite, amsrefs; the wrapper is the lane's pinned \texttt{amsart} editorial wrapper at 10pt on A4, which restates the theorem-like environments the retained text needs and adds the article's own macro definitions that the retained text needs (none), with extra wrapper packages (bm, bbm, dsfont, xspace, cancel, mathtools). The delivered numbering is the unit's own. Assets: no retained span contains a figure environment, a graphics-inclusion command, a PDF-page inclusion, a table or a TikZ picture; no image file is copied into this package, the package declares no asset dependency and no image file is redistributed with it. Licence: version arXiv:2602.17059v2 of this article is distributed under the Creative Commons Attribution 4.0 International licence, as recorded in the arXiv licence metadata for this exact version and shown on the abstract page https://arxiv.org/abs/2602.17059v2. A full rights notice is delivered at licenses/arxiv-2602.17059-CC-BY-4.0.txt. Characters: every retained excerpt is pure ASCII, so the delivered engine, XeLaTeX with its default Latin Modern text font, reports no missing character.
\begin{equation}\label{Linearized_NS}
\left\{\begin{aligned}
\Delta \phi_t + u_s \Delta \phi_x - \Delta u_s \phi_x + v_s \Delta \phi_y - \Delta v_s \phi_y - \varepsilon \Delta^2 \phi &= 0, \quad \mbox{in}~\Omega \times (0,\infty),\\
\phi|_{t=0} &= \phi_0,
\end{aligned}\right.
\end{equation}

\begin{equation}\label{Navier_bc2}
(\phi_y - A\varepsilon^{1/3} \Delta \phi)|_{y=0} = (\phi_y + A\varepsilon^{1/3} \Delta \phi)|_{y=H}= \phi|_{y = 0} = \phi|_{y = H} =0.
\end{equation}\\

\begin{equation}\label{ss_estimates_H}
\begin{aligned}
&0 < u_s(x,y) - u_0(y) \lesssim A \varepsilon^{\frac{1}{3}}, \quad |v_{s}|  \lesssim \varepsilon^{1-}, \quad \|u_{sy} - u_{0y}\|_{L^2} \lesssim \varepsilon^{\frac{2}{3}}, \quad
\|u_{sx}\|_{L^2}  \lesssim \varepsilon^{1-},\\
& |\Delta u_a| \lesssim 1, \qquad | \Delta u_{ax}| + |\Delta u_a - u_{0yy}| \lesssim \varepsilon^{\frac{1}{3}},  \qquad |\Delta v_a| + |\Delta v_{ax}| \lesssim \varepsilon^{\frac{2}{3}},\\
& \| \sqrt{u_0} (\Delta u_s - \Delta u_a) \|_{L^2} + \| \sqrt{u_0} (\Delta v_s - \Delta v_a) \|_{L^2}    \lesssim \varepsilon^{\frac{5}{6}-},
\end{aligned}
\end{equation}

\begin{equation}\label{H1_control_by_R}
\|\phi_x\|_2 + \|\phi_y\|_2 \lesssim \|u_0 \nabla q\|_2 \lesssim H \|R\|_2,
\end{equation}

\begin{corollary}
For $\sigma > 0$, we have
\begin{equation}\label{est:H_weighted_Hardy}
\begin{aligned}
\varepsilon^{\frac{1}{3}}\| \Delta \phi\|_2 & \lesssim \frac{\sqrt\sigma}{\sqrt{H}} \sqrt\varepsilon \| \sqrt{\frac{u_0}{-u_{0yy}}} \Delta \phi_y\|_2 + \sigma^{-1} \|\frac{R}{\sqrt{-u_{0yy}}}\|_2,\\
\varepsilon^{\frac{1}{6}}\| \sqrt{\frac{u_0}{-u_{0yy}}}  \Delta \phi\|_2 & \lesssim \sigma \sqrt\varepsilon \| \sqrt{\frac{u_0}{-u_{0yy}}} \Delta \phi_y\|_2 + \sigma^{-1/2} \sqrt{H}  \|\frac{R}{\sqrt{-u_{0yy}}}\|_2.\\
\end{aligned}
\end{equation}
\end{corollary}

\begin{equation}\label{est:mix_norm_H}
\begin{aligned}
\| \nabla \phi\|_{L^2_x L^\infty_y} &\lesssim_H \varepsilon^{-\frac{1}{6}} \Big( \|\frac{R}{\sqrt{-u_{0yy}}}\|_2 + \sqrt{\varepsilon} \| \sqrt{\frac{u_0}{-u_{0yy}}} \Delta \phi_y\|_2\Big),\\
\| \nabla \phi\|_{L^\infty_x L^2_y} &\lesssim_H \varepsilon^{-\frac{1}{6}}  \Big( \|\frac{R}{\sqrt{-u_{0yy}}}\|_2 + \sqrt{\varepsilon} \| \sqrt{\frac{u_0}{-u_{0yy}}} \Delta \phi_{y}\|_2\Big).
\end{aligned}
\end{equation}

\begin{equation}\label{est:mix_norm_Hardy}
\begin{aligned}
\| \frac{\phi}{\sqrt{u_0}}\|_{L^2_x L^\infty_y} &\lesssim_H \|\phi_y\|_2 \lesssim_H \|\frac{R}{\sqrt{-u_{0yy}}}\|_2,\\
\| \frac{\phi_x}{\sqrt{u_0}}\|_{L^2_x L^\infty_y} &\lesssim_H \|\phi_{xy}\|_2 \lesssim_H \varepsilon^{-\frac{1}{3}} \Big( \|\frac{R}{\sqrt{-u_{0yy}}}\|_2 + \sqrt{\varepsilon} \| \sqrt{\frac{u_0}{-u_{0yy}}} \Delta \phi_{y}\|_2\Big).
\end{aligned}
\end{equation}

\begin{lemma}
\begin{equation}\label{est:mix_norm_misc}
\begin{aligned}
\| \sqrt{u_0} \Delta \phi\|_{L^\infty_x L^2_y} &\lesssim_{H} \varepsilon^{-\frac{1}{3}} \Big( \|\frac{R}{\sqrt{-u_{0yy}}}\|_2 + \sqrt{\varepsilon} \| \sqrt{\frac{u_0}{-u_{0yy}}} \nabla \Delta \phi\|_2\Big),\\
\|\phi\|_{L^\infty_xL^2_y} &\lesssim_{L}  \|\frac{R}{\sqrt{-u_{0yy}}}\|_2.
\end{aligned}
\end{equation}
\end{lemma}

\begin{equation}\label{est:g_basic}
\|g\|_2 + \| \nabla g\|_2 + \| \nabla^2 g\|_2 \lesssim_{H,L} \| \phi_t\|_2.
\end{equation}

\begin{lemma}
\begin{equation}\label{est:g_x}
\|g_x\|_2 \le \frac{H}{4\pi} \|\phi_t\|_2.
\end{equation}
\end{lemma}

\begin{proposition}\label{prop:phi_t}
For any $\delta > 0$, we can choose an $N$ large from the weight $W = NL - x$, and there exists some positive constant $C$ depending only on $\delta, A, L$, and $H$, such that
\begin{equation}\label{est:phi_t}
\|\phi_t\|_2 \le \frac{H}{4\pi}\|R\|_2 + \frac{\delta}{\sqrt{2}} \| \sqrt{\frac{W}{-u_{0yy}}} R\|_{x= L} + C \varepsilon^{\frac{1}{3}}  \Big( \|\frac{R}{\sqrt{-u_{0yy}}}\|_2 + \sqrt{\varepsilon} \| \sqrt{\frac{u_0}{-u_{0yy}}} \nabla \Delta \phi\|_2\Big).
\end{equation}
\end{proposition}

\begin{proof}
We rewrite the linearized Navier-Stokes equation \eqref{Linearized_NS} as
$$
\Delta \phi_t = - R_x -(u_s - u_0) \Delta \phi_x + (\Delta u_s - u_{0yy}) \phi_x - v_s \Delta \phi_y + \Delta v_s \phi_y + \varepsilon \Delta^2 \phi.
$$
We test the equation by $-g$ and estimate term by term.

For the temporal term, we simply integrate by parts and obtain
\begin{equation}\label{Dual_1}
(\Delta \phi_t, -g) = (\phi_t , - \Delta g) = \| \phi_t\|_2^2.
\end{equation}

For the Rayleigh terms, first by Sobolev inequality in $x$ and \eqref{est:g_x}, we have
\begin{align*}
|(R_x, g)| &\le |-(R,g_x)| + |(R,g)_{x=L}|\\
&\le \|R\|_2 \|g_x\|_2 + \sqrt{L}\|R\|_{x=L}\|g_x\|_2\\
&\le \frac{H}{4\pi} \|R\|_2 \|\phi_t\|_2 + \frac{H\sqrt{L}}{4\pi} \|R\|_{x=L} \| \phi_t\|_2.
\end{align*}
Then for any $\delta > 0$, we can choose $N$ large such that
$$
\frac{\sqrt{2}H}{4\pi \sqrt{-(N-1)u_{0yy}}} < \delta,
$$
which implies
\begin{equation}\label{Dual_2}
|(R_x, g)| \le \frac{H}{4\pi} \|R\|_2 \|\phi_t\|_2 + \frac{\delta}{\sqrt{2}} \| \sqrt{\frac{W}{-u_{0yy}}} R\|_{x= L} \| \phi_t\|_2.
\end{equation}

For the next Rayleigh term, we integrate by parts and write
\begin{equation}\label{Dual_3}
((u_s - u_0) \Delta \phi_x,g) = - ((u_s - u_0) \Delta \phi, g_x) + ((u_s- u_0) \Delta \phi ,g)_{x=L} - (u_{sx} \Delta \phi, g).
\end{equation}
For the first two terms, we use \eqref{ss_estimates_H}, \eqref{H1_control_by_R}, Hardy inequality in $y$, Sobolev inequality in $x$, and \eqref{est:g_basic} to estimate
\begin{equation}\label{Dual_3-1}
\begin{aligned}
|(\ref{Dual_3}.1+2)| \lesssim& A \varepsilon^{\frac{1}{3}} \| u_0 \Delta \phi\|_2 \| \frac{g_x}{u_0}\|_2 + A \varepsilon^{\frac{1}{3}} \| u_0 \Delta \phi\|_{x=L} \| \frac{g}{u_0}\|_{x=L}\\
\lesssim&_{A,H,L}~ \varepsilon^{\frac{1}{3}} \| \frac{R}{\sqrt{-u_{0yy}}} \|_2 \| g_{xy}\|_2 + \varepsilon^{\frac{1}{3}} \Big( \| \frac{R}{\sqrt{-u_{0yy}}} \|_{x=L} + \| \phi\|_{x=L} \Big) \| g_y\|_{x=L}\\
\lesssim&_{A,H,L}~ \varepsilon^{\frac{1}{3}} \| \frac{R}{\sqrt{-u_{0yy}}} \|_2 \| g_{xy}\|_2 + \varepsilon^{\frac{1}{3}} \Big( \| \frac{R}{\sqrt{-u_{0yy}}} \|_{x=L} + \| \phi_x\|_{2} \Big) (\| g_y\|_{2} + \|g_{xy}\|_2)\\
\lesssim&_{A,H,L}~ \varepsilon^{\frac{1}{3}} \Big( \| \frac{R}{\sqrt{-u_{0yy}}} \|_2 +  \| \frac{R}{\sqrt{-u_{0yy}}} \|_{x=L} \Big) \| \phi_t\|_2.
\end{aligned}
\end{equation}
For the third term, we use \eqref{ss_estimates_H}, \eqref{est:mix_norm_misc}, \eqref{est:mix_norm_Hardy} with $\phi = g$, and \eqref{est:g_basic} to estimate
\begin{equation}\label{Dual_3-2}
\begin{aligned}
|(\ref{Dual_3}.3)| \lesssim& \|u_{sx}\|_2 \| \sqrt{u_0} \Delta \phi\|_{L^\infty_x L^2_y} \| \frac{g}{\sqrt{u_0}}\|_{L^2_x L^\infty_y}\\
\lesssim&_{A,H,L} \varepsilon^{1-} \| \sqrt{u_0} \Delta \phi\|_{L^\infty_x L^2_y}  \|g_y\|_2\\
\lesssim&_{A,H,L} \varepsilon^{\frac{2}{3}-} \Big( \|\frac{R}{\sqrt{-u_{0yy}}}\|_2 + \sqrt{\varepsilon} \| \sqrt{\frac{u_0}{-u_{0yy}}} \nabla \Delta \phi\|_2\Big) \|\phi_t\|_2.
\end{aligned}
\end{equation}

Next, we split $\Delta u_s$ and estimate using \eqref{ss_estimates_H}, \eqref{est:mix_norm_misc}, \eqref{est:mix_norm_H}, \eqref{est:mix_norm_Hardy} with $\phi = g$, and \eqref{est:g_basic}
\begin{equation}\label{Dual_4}
\begin{aligned}
|((\Delta u_s - u_{0yy}) \phi_x, g)| \le& |((\Delta u_s - \Delta u_a) \phi_x, g) | + |((\Delta u_a - u_{0yy}) \phi_x, g) |\\
\lesssim&_{A,H,L} \| \sqrt{u_0} (\Delta u_s - \Delta u_a)\|_2 \| \phi_x\|_{L^\infty_x L^2_y} \| \frac{g}{\sqrt{u_0}}\|_{L^2_x L^\infty_y} + \varepsilon^{\frac{1}{3}}\|\phi_x\|_2\|g\|_2\\
\lesssim&_{A,H,L}~ \varepsilon^{\frac{1}{3}}  \Big( \|\frac{R}{\sqrt{-u_{0yy}}}\|_2 + \sqrt{\varepsilon} \| \sqrt{\frac{u_0}{-u_{0yy}}} \nabla \Delta \phi\|_2\Big) \|\phi_t\|_2.
\end{aligned}
\end{equation}

Similarly, we can split $\Delta v_s$ and estimate
\begin{equation}\label{Dual_5}
\begin{aligned}
|(\Delta v_s \phi_y, g)| \le& |((\Delta v_s - \Delta v_a) \phi_y, g) | + |(\Delta v_a  \phi_x, g) |\\
\lesssim&_{A,H,L} \| \sqrt{u_0} (\Delta v_s - \Delta v_a)\|_2 \| \phi_y\|_{L^\infty_x L^2_y} \| \frac{g}{\sqrt{u_0}}\|_{L^2_x L^\infty_y} + \varepsilon^{\frac{2}{3}}\|\phi_y\|_2\|g\|_2\\
\lesssim&_{A,H,L}~ \varepsilon^{\frac{2}{3}-}  \Big( \|\frac{R}{\sqrt{-u_{0yy}}}\|_2 + \sqrt{\varepsilon} \| \sqrt{\frac{u_0}{-u_{0yy}}} \nabla \Delta \phi\|_2\Big) \|\phi_t\|_2.
\end{aligned}
\end{equation}

For the last Rayleigh term, we can simply use \eqref{ss_estimates_H}, Hardy inequality in $y$, and \eqref{est:g_basic} to estimate
\begin{equation}\label{Dual_6}
\begin{aligned}
|( v_s \Delta\phi_y, g)| \lesssim_{A,H,L} \varepsilon^{1-} \| \sqrt{u_0} \Delta \phi_y\|_2\| \frac{g}{\sqrt{u_0}}\|_2 \lesssim_{A,H,L} \varepsilon^{\frac{1}{2}-} \Big( \sqrt{\varepsilon} \| \sqrt{\frac{u_0}{-u_{0yy}}}  \Delta \phi_y\|_2\Big)\phi_t\|_2.
\end{aligned}
\end{equation}

Finally for the dissipation term, we integrate by parts and use the $\varepsilon$-Navier boundary condition \eqref{Navier_bc2} to obtain
\begin{equation*}
\begin{aligned}
-\varepsilon(\Delta^2 \phi, g) &= - \varepsilon (\Delta \phi, \Delta g) + \varepsilon (\Delta \phi, g_y)_{y = H} - \varepsilon (\Delta \phi, g_y)_{y = 0}\\
&= - \varepsilon (\Delta \phi, \Delta g) - \frac{\varepsilon^{\frac{2}{3}}}{A} ( \phi_y, g_y)_{y = H} - \frac{\varepsilon^{\frac{2}{3}}}{A} ( \phi_y, g_y)_{y = 0}.
\end{aligned}
\end{equation*}
Therefore, by Sobolev inequality in $y$, \eqref{H1_control_by_R}, and \eqref{est:H_weighted_Hardy}, we can estimate
\begin{equation}\label{Dual_7}
\begin{aligned}
\varepsilon|(\Delta^2 \phi, g)| &\lesssim_{A,H} \varepsilon \| \Delta\phi\|_2\|\Delta g\|_2 + \varepsilon^{\frac{2}{3}} (\| \phi_y\|_2 + \| \phi_{yy}\|_2) (\|g_y\|_2 + \| g_{yy}\|_2)\\
&\lesssim_{A,H} \varepsilon^{\frac{1}{3}}  \Big( \|\frac{R}{\sqrt{-u_{0yy}}}\|_2 + \sqrt{\varepsilon} \| \sqrt{\frac{u_0}{-u_{0yy}}} \nabla \Delta \phi\|_2\Big) \|\phi_t\|_2.
\end{aligned}
\end{equation}
Collecting all the terms from \eqref{Dual_1}-\eqref{Dual_7} yields the desired estimate \eqref{est:phi_t}. This concludes the proof.
\end{proof}

\end{document}
